\chapter{Antywzorce i pulapki}

\section{Po co poznawac antywzorce}
Dobre praktyki mowia, co robic. Antywzorce pokazuja, czego unikac. To równie wazne,
bo wiele problemow nie wynika z braku wiedzy, tylko z pozornie wygodnych skrótow,
ktore w dluzszej perspektywie psuja czytelnosc albo stabilnosc projektu.

\section{Antywzorzec: wszystko w jednym pliku}
Poczatkujacy często zaczynaja od jednego ogromnego pliku i zostaja przy nim zbyt dlugo.
Skutek jest prosty: metadata, funkcje, dane, strony i debug printy mieszaja sie ze soba.

Jesli projekt zaczyna rosnac, podziel go na chunki i pliki. Nie czekaj, az zrobi sie
naprawde nieczytelny.

\section{Antywzorzec: kopiowanie metadata}
Kopiowanie tych samych pol \mc{title}, \mc{metaDescription}, \mc{lang} czy \mc{styles}
do wielu chunkow predzej czy pozniej doprowadzi do rozjazdow. Jedna strona bedzie miala
stary opis, inna nowy, trzecia w ogole bez aktualizacji.

Lepiej wydzielic wspolny chunk \mc{Metadata} i uzywac \mc{@include}.

\section{Antywzorzec: zbyt duzo logiki w \mc{content}}
To jeden z najczestszych bledow. Jesli w srodku HTML-a masz kilka zagniezdzonych warunkow,
dlugie rzutowania i budowanie tekstu na piec sposobow, to znak, ze logika powinna zostac
przeniesiona wyzej do zmiennych albo funkcji.

\section{Antywzorzec: naduzywanie \mc{unknown} i \mc{any}}
Te typy sa potrzebne, ale nie powinny stac sie codziennym domyslnym rozwiazaniem.
Jesli wszedzie uciekasz do \mc{unknown} albo \mc{any}, tak naprawde rezygnujesz z duzej
części korzysci wynikajacych ze swiadomego modelu typow.

\section{Antywzorzec: brak warunku stopu w rekurencji}
To klasyczna pulapka. Rekurencja bez sensownego warunku stopu nie prowadzi do wyniku,
tylko do spirali wywolan. Jesli piszesz funkcje wywolujaca sama siebie, pierwszym pytaniem
powinno byc zawsze: \emph{kiedy to sie zatrzyma?}

\section{Antywzorzec: nieczytelne nazwy}
Nazwy typu \mc{tmp2}, \mc{val}, \mc{x1} i \mc{something} nie pomagaja zrozumiec kodu.
Zwlaszcza w DSL-u do generowania stron dobra nazwa potrafi zaoszczedzic bardzo duzo czasu.

\section{Antywzorzec: traktowanie \mc{if} jako miejsca do swobodnego cieniowania}
MoonChunk celowo jest rygorystyczny w kwestii redeklaracji nazw w cialach \mc{if},
\mc{for} i \mc{while}. Jesli sprobujesz tworzyc tam nowe zmienne o tych samych nazwach,
które juz istnieja w poblizu, szybko trafisz na komunikat o redeklaracji.

Jesli naprawde potrzebujesz osobnego zakresu z ta sama nazwa, uzyj jawnego bloku
\mc{\{ ... \}} albo funkcji.

\section{Antywzorzec: import bez realnego uzycia}
Samo importowanie pliku niczego nie uruchamia. Jesli liczysz na to, ze zaimportowany
chunk wykona sie sam, projekt zacznie wygladac tak, jakby ``nic nie dzialalo''.
Musisz go swiadomie dolaczyc przez \mc{@include} albo uruchomic przez \mc{moon(...)}.

\section{Antywzorzec: brak strategii debugowania}
Chaotyczne dodawanie losowych \mc{print(...)} jest lepsze niz nic, ale nie duzo lepsze.
Lepsza strategia to zadac sobie pytanie: jaka wartosc chce sprawdzic i w ktorym miejscu?
Dopiero wtedy dodaj log.

\section{Antywzorzec: projekt budowany od konca}
Czesta pomylka polega na tym, ze autor zaczyna od bardzo rozbudowanej struktury, zanim
powstanie pierwszy dzialajacy przebieg. Lepiej zrobic odwrotnie: najpierw maly projekt,
ktory naprawdę generuje HTML, a dopiero potem dokladac warstwy i modularyzacje.
